\subsection{Sealed Bids Method}

\content{
        \begin{definition} The method begins by compiling a list of items to be
        divided. Then: 
        \begin{enumerate}
        \item Each party involved lists, in secret, a dollar amount they 
        value each item to be worth. This is their sealed bid.
        \item The bids are collected. For each party, the value of all the 
        items is totaled, and divided by the number of parties. This 
        defines their fair share.  
        \item Each item is awarded to the highest bidder.
        \item For each party, the value of all items received is totaled. 
        If the value is more than that party’s fair share, they pay the 
        difference into a holding pile. If the value is less than that party
        ’s fair share, they receive the difference from the holding pile. 
        This ends the initial allocation.  
        \item In most cases, there will be a surplus, or leftover money, in 
        the holding pile.  The surplus is divided evenly between all the 
        players. This produces the final allocation.
        \end{enumerate}
        \end{definition}
}{%presentation
}{% nopresentation
}{% comments
}

\content{
        \begin{example} Sam and Omar have cohabitated for the last 3 years, 
        during which time they shared the expense of purchasing several 
        items for their home. Sam has accepted a job in another city, and 
        now they find themselves needing to divide their shared assets.  
        Each records their bids for each item as shown below. \\

        \begin{tabular}{|l|l|l|}
        \hline
        & Sam & Omar \\
        \hline
        Couch & \$150 & \$100 \\
        \hline
        TV & \$200 & \$250 \\
        \hline
        Game Console & \$250 & \$150 \\
        \hline
        Sound System & \$50 & \$100\\
        \hline
        \end{tabular}
        \end{example}
}{%presentation
\vspace{3in}
}{% nopresentation
\vspace{3in}\newpage
}{% comments
}

\content{
        \begin{example} Jamal, Maggie, and Kendra are dividing an estate 
        consisting of a house, a vacation home, and a small business. Their 
        valuations (in thousands) are shown below. Determine the final 
        allocation.  \\

        \begin{tabular}{|l|l|l|l|}
        \hline
        & Jamal & Maggie & Kendra \\
        \hline
        House & \$250 & \$300 & \$280 \\
        \hline
        Vacation Home & \$170 & \$180 & \$200 \\
        \hline
        Small Business & \$300 & \$255 & \$270 \\
        \hline
        \end{tabular}
        \end{example}
}{%presentation
\vspace{3in}
}{% nopresentation
\vspace{4in}
}{% comments
}

\content{
        \begin{example} As part of an inheritance, four children, Abby, Ben 
        and Carla, are dividing four vehicles using Sealed Bids. Their bids 
        (in thousands of dollars) for each item is shown below. Find the 
        final allocation.\\

        \begin{tabular}{|l|l|l|l|}
        \hline
        & Abby & Ben & Carla \\
        \hline
        Motorcycle & 10  & 9 & 8 \\
        \hline
        Car & 10 & 11 & 9 \\
        \hline
        Tractor & 4 & 1 & 2 \\
        \hline
        Boat & 7 & 6 & 4 \\
        \hline
        \end{tabular}
        \end{example}
}{%presentation
\vspace{3in}
}{% nopresentation
\vspace{3in}
}{% comments
}

\content{
        \begin{example} Fair division does not always have to be used for 
        items of value. It can also be used to divide undesirable items. 
        Suppose Chelsea and Mariah are sharing an apartment, and need to 
        split the chores for the household. They list the chores, assigning 
        a negative dollar value to each item; in other words, the amount 
        they would pay for someone else to do the chore (a per week amount).
        We will assume, however, that they are committed to doing all 
        the chores themselves and not hiring a maid. Find a fair division if the
        bids are given as shown below. \\

        \begin{tabular}{|l|l|l|}
        \hline
        & Chelsea & Mariah \\
        \hline
        Vacuuming & -\$10 & -\$8 \\
        \hline
        Cleaning bathroom  & -\$14 & -\$20 \\
        \hline
        Doing Dishes & -\$4 & -\$6 \\
        \hline
        Dusting & -\$6 & -\$4 \\
        \hline
        \end{tabular}
        \end{example} 
}{%presentation
\vspace{3in}
}{% nopresentation
\vspace{3in}
}{% comments
}

